\subsection{一点处的芽}

第 9 目的定义与结果所适用的最常见的情形之一是：\(\mathfrak{F}\) 是一点 \(a\)（拓扑空间 \(\mathbf{X}\) 的）的邻域滤子；这时我们不谈“关于 \(\mathfrak{F}\) 的芽”，而谈“在点 \(a\) 处的芽”。注意，\(a\) 的邻域只有一个芽，即整个空间 \(\mathbf{X}\) 的芽。闭集的芽与在点 \(a\) 处局部闭的集合的芽相同，因为若 \(\mathbf{L}\) 在 \(a\) 处局部闭，则 \(\mathbf{L}\) 与 \(\overline{\mathbf{L}}\) 在 \(a\) 处的芽相等（\(\S 3\)，第 1 目，命题 1）。由此可知，若 \(\xi, \eta\) 是在 \(a\) 处局部闭的集合的两个芽，则 \(\xi \cup \eta\) 与 \(\xi \cap \eta\) 也是这样的芽。

由于 \(a\) 属于每个 \(\mathbf{V} \in \mathfrak{F}\)，故 \(f(a)\) 对每个映射 \(f\)（定义域属于 \(\mathfrak{F}\) 的）都有定义；此外，若 \(f\) 与 \(g\) 在 \(a\) 处有相同的芽，则必有 \(f(a) = g(a)\)，于是 \(f(a)\) 只依赖于芽 \(\tilde{f}\)（\(f\) 在 \(a\) 处的），称为 \(\tilde{f}\) 在 \(a\) 的值，记作 \(\tilde{f}(a)\)。应当强调，关系 \(\tilde{f}(a) = \tilde{g}(a)\) 一般并不蕴含 \(\tilde{f} = \tilde{g}\)。

设 \(X', X''\) 是两个拓扑空间；\(b\) 是 \(X''\) 的一点；\(g, g'\) 是 \(X'\) 到 \(X''\) 中的两个映射，它们在 \(b\) 处有相同的芽。若 \(f, f'\) 是 \(X\) 到 \(X'\) 中的两个映射，都在 \(a\) 处连续、在 \(a\) 处有相同的芽，且满足 \(f(a) = b\)，则 \(g \circ f\) 与 \(g' \circ f'\) 在点 \(a\) 处有相同的芽；因为若 \(V'\) 是 \(b\) 在 \(X'\) 中的一个邻域，使得 \(g(x') = g'(x')\)（对所有 \(x' \in V'\)），则存在邻域 \(V\)（\(a\) 的），使得 \(f(V) \subset V'\)、\(f'(V) \subset V'\)，且 \(f(x) = f'(x)\)（对所有 \(x \in V\)），从而断言得证。\(g \circ f\) 在 \(a\) 处的芽称为芽 \(\tilde{g}\) 与 \(\tilde{f}\)（\(g\) 与 \(f\) 的）的合成，记作 \(\tilde{g} \circ \tilde{f}\)。

